% Ý nghĩa của cận trên LP và gap (minh họa).
\documentclass[border=6pt]{standalone}
% Shared preamble of the standalone figures used by docs/explaination.md.
% Vietnamese text needs T5 + babel (the report itself is English, T1).
\usepackage[utf8]{inputenc}
\usepackage[T1,T5]{fontenc}
\usepackage[english,vietnamese]{babel}
\usepackage{lmodern}
\usepackage{amsmath,amssymb}
\usepackage{xcolor}
\usepackage{tikz}
\usepackage{pgfplots}
\pgfplotsset{compat=1.18}
\usepackage{fontawesome5}
\definecolor{groundgreen}{RGB}{168,196,116}
\usetikzlibrary{arrows.meta, positioning, calc, fit, backgrounds, shapes.geometric, shapes.misc, decorations.pathmorphing, decorations.pathreplacing, patterns}
\newcommand{\R}{\mathbb{R}}
\newcommand{\norm}[1]{\lVert #1\rVert}
\renewcommand{\familydefault}{\sfdefault}

\begin{document}
\begin{tikzpicture}[x=0.38cm,y=1cm,>={Stealth[length=2.2mm]}]
\draw[->, thick, black!70] (0,0) -- (42,0) node[right, font=\small] {số alert đến kịp};
\foreach \t in {0,10,20,30,40} \draw (\t,0.12) -- (\t,-0.12) node[below, font=\small] {\t};
% achieved
\fill[blue!60!black] (28,0) circle (5pt);
\node[font=\small, align=center, text=blue!60!black, anchor=north] at (28,-0.55) {SBS đạt được\\(28 alert)};
% bound
\fill[red!70!black] (35,0) circle (5pt);
\node[font=\small, align=center, text=red!70!black, anchor=north] at (35,-0.55) {cận trên LP\\(35 alert)};
% unknown optimum zone
\fill[orange!25] (28.4,0.25) rectangle (34.6,0.9);
\node[font=\small, align=center] at (31.5,0.57) {tối ưu thật\\nằm đâu đó ở đây};
\draw[decorate, decoration={brace, amplitude=6pt}, thick] (28,1.15) -- node[above=7pt, font=\small, align=center] {gap $= \dfrac{35-28}{35} = 0{,}20$\\ = \emph{trần} của phần còn cải thiện được} (35,1.15);
\node[font=\small, align=left, anchor=north west, text width=10cm] at (0,2.6) {%
  \textbf{Cận trên} = kết quả của một bài toán \emph{nới lỏng} (LP):\\
  -- cho phép chọn \emph{nửa đường} ($0 \le x \le 1$)\\
  -- thay đường cong tốc độ bằng 10 tiếp tuyến nằm trên nó\\
  $\Rightarrow$ dễ hơn và rộng hơn, nên \textbf{không kế hoạch nào vượt qua được}.};
\node[font=\scriptsize, text=black!60, anchor=west] at (0,-1.8) {(minh họa; số thật ở Table 5: gap của SBS = 0,191, FixedTour = 0,238)};
\end{tikzpicture}
\end{document}
